\section*{第 \ref{ch:g7:what-is-respiration} 章——什么是呼吸作用？}

\begin{solution}{exo:g7:what-is-respiration:1}
\omterm{def:g7:what-is-respiration:respiration}{呼吸作用}是一种气体交换，\omterm{def:g6:exploring-our-environment:components}{生物}通过它从
外界获取\omterm{def:g7:what-is-respiration:respiration}{氧气}，并释放\omterm{def:g7:what-is-respiration:respiration}{二氧化碳}。
\omterm{def:g7:what-is-respiration:respiration}{呼吸}只是胸部可见的一起一伏的运动——
它只是服务于这种交换的一种方式；交换本身发生在
每一个活\omterm{def:g5:first-look-at-cells:cell}{细胞}中。
\end{solution}

\begin{solution}{exo:g7:what-is-respiration:2}
密闭容器中的\omterm{def:g7:what-is-respiration:respiration}{氧气}探头——显示\omterm{def:g7:what-is-respiration:respiration}{氧气}
含量下降；以及\omterm{prop:g7:what-is-respiration:shown}{澄清石灰水}——一种透明的液体，
\omterm{def:g7:what-is-respiration:respiration}{二氧化碳}通入其中时会变浑浊。
\end{solution}

\begin{solution}{exo:g7:what-is-respiration:3}
为了证明这一变化来自活的研究对象，而不是来自
容器、白天的\omterm{def:g6:exploring-our-environment:conditions}{温度}或\omterm{prop:g7:what-is-respiration:shown}{澄清石灰水}本身，
这正是公平实验中的对照，唯一的差别在于有没有生命。
\end{solution}

\begin{solution}{exo:g7:what-is-respiration:4}
\omterm{prop:g3:balanced-diet:jobs}{能量}：在每个\omterm{def:g5:first-look-at-cells:cell}{细胞}中，\omterm{def:g4:journey-of-food:digestion}{营养物质}与\omterm{def:g7:what-is-respiration:respiration}{氧气}发生缓慢的
燃烧，释放出\omterm{prop:g3:balanced-diet:jobs}{能量}用于\omterm{def:g5:first-look-at-cells:cell}{细胞}的工作。\omterm{def:g7:what-is-respiration:respiration}{二氧化碳}
则是燃烧后的剩余物，被运走。
\end{solution}

\begin{solution}{exo:g7:what-is-respiration:5}
一只小鼠、一只潮虫、正在\omterm{def:g2:from-seed-to-plant:germination}{萌发}的豌豆、一朵蘑菇、池塘
螺——以及一勺活的\omterm{def:g6:decomposers-and-soil:soil}{土壤}。
\end{solution}

\begin{solution}{exo:g7:what-is-respiration:6}
在黑暗中检测：光驱动的物质制造停止后，
广口瓶中的\omterm{def:g7:what-is-respiration:respiration}{氧气}减少，\omterm{prop:g7:what-is-respiration:shown}{澄清石灰水}变浑浊——
\omterm{def:g7:what-is-respiration:respiration}{呼吸作用}便显现出来。
\end{solution}

\begin{solution}{exo:g7:what-is-respiration:7}
相同点：两者都消耗\omterm{def:g7:what-is-respiration:respiration}{氧气}，释放\omterm{def:g7:what-is-respiration:respiration}{二氧化碳}；
两者都释放燃料中的\omterm{prop:g3:balanced-diet:jobs}{能量}。不同点：\omterm{def:g5:first-look-at-cells:cell}{细胞}内的
燃烧缓慢且无火焰，分成许多小步骤进行；
其\omterm{prop:g3:balanced-diet:jobs}{能量}被储存起来用于做功，而不是
一下子以光和热的形式散失。
\end{solution}

\begin{solution}{exo:g7:what-is-respiration:8}
正在冲刺的儿童（剧烈的\omterm{def:g3:bones-and-muscles:muscle}{肌肉}运动）$>$ 正在阅读的儿童（
身体安静时的基础消耗）$>$ 正在\omterm{def:g2:from-seed-to-plant:germination}{萌发}的豌豆（生长迅速，
但体积小）$>$ 干燥的豌豆（休眠——\omterm{def:g7:what-is-respiration:respiration}{呼吸作用}
几乎停滞）。速率随做功而变化。
\end{solution}

\begin{solution}{exo:g7:what-is-respiration:9}
整夜，\omterm{def:g2:from-seed-to-plant:germination}{萌发}的豌豆进行\omterm{def:g7:what-is-respiration:respiration}{呼吸作用}，消耗了
广口瓶中的\omterm{def:g7:what-is-respiration:respiration}{氧气}，并释放出\omterm{def:g7:what-is-respiration:respiration}{二氧化碳}；到了早晨，
剩余的\omterm{def:g7:what-is-respiration:respiration}{氧气}已不足以维持火焰燃烧。\omterm{prop:g7:what-is-respiration:shown}{澄清石灰水}
整夜都会持续变浑浊。
\end{solution}

\begin{solution}{exo:g7:what-is-respiration:10}
两个装有探头和\omterm{prop:g7:what-is-respiration:shown}{澄清石灰水}吸收装置的容器：一个装
新鲜的活\omterm{def:g6:decomposers-and-soil:soil}{土壤}，另一个装同样的\omterm{def:g6:decomposers-and-soil:soil}{土壤}，但经烘烤失去生命
——唯一的差别在此。若活\omterm{def:g6:decomposers-and-soil:soil}{土壤}的广口瓶中\omterm{def:g7:what-is-respiration:respiration}{氧气}下降，
\omterm{prop:g7:what-is-respiration:shown}{澄清石灰水}变浑浊，而经烘烤的对照瓶保持不变，
则说明是\omterm{def:g6:decomposers-and-soil:soil}{土壤}中的生命所致。
\end{solution}

\begin{solution}{exo:g7:what-is-respiration:11}
潮湿的谷粒苏醒：\omterm{def:g2:from-seed-to-plant:germination}{萌发}开始，\omterm{def:g7:what-is-respiration:respiration}{呼吸作用}由微弱
逐渐升至接近\omterm{def:g2:from-seed-to-plant:germination}{萌发}豌豆的速率——数十亿粒谷粒
同时消耗各自储备的物质，使谷堆发热，
就像堆肥中忙碌的\omterm{def:g6:producing-our-food:fermentation}{微生物}群一样。\omterm{def:g6:exploring-our-environment:conditions}{温度}升高会加快这一过程，
进一步促使潮湿的谷堆自行升温，可能导致变质。干燥的
谷粒处于低温休眠；潮湿的谷粒则像一座缓慢燃烧的炉子。
\end{solution}

\begin{solution}{exo:g7:what-is-respiration:12}
行为：\omterm{def:g7:what-is-respiration:respiration}{呼吸}可以被观察到——胸部的
起伏运动，每分钟十五次
（见\cref{ex:g4:breathing:watch}）。化学：\omterm{def:g7:what-is-respiration:respiration}{呼吸作用}发生在
完全没有胸部的\omterm{def:g6:exploring-our-environment:components}{生物}中——豌豆、蘑菇、\omterm{def:g6:decomposers-and-soil:soil}{土壤}
——只能通过气体检测发现。功能：\omterm{def:g7:what-is-respiration:respiration}{呼吸}的存在
是为了更新肺内的空气，使气体交换得以持续——
其加快恰与\omterm{def:g5:first-look-at-cells:cell}{细胞}燃烧的加快同步
（见\cref{ex:g4:breathing:running}）。
\end{solution}

\begin{solution}{pb:g7:what-is-respiration:1}
\textbf{1.} 无变化：\omterm{def:g7:what-is-respiration:respiration}{氧气}含量稳定，\omterm{prop:g7:what-is-respiration:shown}{澄清石灰水}仍然澄清。
A 是无生命的对照，
其余各瓶都以它的稳定状态为参照。

\textbf{2.} B：几乎测不到变化——干燥的豌豆处于
休眠，\omterm{def:g7:what-is-respiration:respiration}{呼吸作用}几乎停滞。C：\omterm{def:g7:what-is-respiration:respiration}{氧气}明显下降并
变浑浊——\omterm{def:g2:from-seed-to-plant:germination}{萌发}正处于全速生长，做功
所需\omterm{prop:g3:balanced-diet:jobs}{能量}由\omterm{def:g7:what-is-respiration:respiration}{呼吸作用}提供。

\textbf{3.} D：下降和变浑浊最快的一瓶——一只
温暖、活动的\omterm{prop:g3:sorting-living-things:groups}{哺乳动物}耗氧很快。实验的时长有限，
因为小鼠正在消耗广口瓶中的\omterm{def:g7:what-is-respiration:respiration}{氧气}；
实验必须远在空气不足之前结束。

\textbf{4.} E：\omterm{def:g7:what-is-respiration:respiration}{氧气}含量稳定或上升——受光照的\omterm{def:g1:plants-around-us:plant}{植物}的
物质制造吸收\omterm{def:g7:what-is-respiration:respiration}{二氧化碳}，释放的\omterm{def:g7:what-is-respiration:respiration}{氧气}
多于其\omterm{def:g7:what-is-respiration:respiration}{呼吸作用}所消耗的。F：\omterm{def:g7:what-is-respiration:respiration}{氧气}下降，
\omterm{prop:g7:what-is-respiration:shown}{澄清石灰水}变浑浊——在黑暗中，物质制造停止，
而从未停止的\omterm{def:g7:what-is-respiration:respiration}{呼吸作用}便清楚地显现出来。

\textbf{5.} F 表明，\omterm{def:g1:plants-around-us:plant}{植物}始终在进行\omterm{def:g7:what-is-respiration:respiration}{呼吸作用}——
吸入\omterm{def:g7:what-is-respiration:respiration}{氧气}，呼出\omterm{def:g7:what-is-respiration:respiration}{二氧化碳}，与小鼠一样。E 则在白天
规模更大的物质制造中掩盖了同样的\omterm{def:g7:what-is-respiration:respiration}{呼吸作用}：
制造过程的吸入与释放量淹没了它。两瓶合在一起，
将两种气体交换分开。

\textbf{6.} 正在\omterm{def:g2:from-seed-to-plant:germination}{萌发}的豌豆，每克的工作强度
是三倍之多——根与芽正以全速生长——
\omterm{def:g7:what-is-respiration:respiration}{呼吸作用}正是这份工作的“引擎声”。

\textbf{7.} 小鼠的\omterm{def:g3:bones-and-muscles:muscle}{肌肉}和温热的身体借助\omterm{def:g7:what-is-respiration:respiration}{氧气}
使\omterm{def:g4:journey-of-food:digestion}{营养物质}燃烧；燃烧释放出的\omterm{def:g7:what-is-respiration:respiration}{二氧化碳}
进入广口瓶内的空气；被困住的空气
通过\omterm{prop:g7:what-is-respiration:shown}{澄清石灰水}，\omterm{def:g7:what-is-respiration:respiration}{二氧化碳}使其变浑浊。

\textbf{8.} “每种\omterm{def:g6:exploring-our-environment:components}{生物}都进行\omterm{def:g7:what-is-respiration:respiration}{呼吸作用}——\omterm{def:g1:animals-around-us:animal}{动物}、
\omterm{def:g1:plants-around-us:plant}{植物}、清醒或几乎睡着的\omterm{def:g2:from-seed-to-plant:seed}{种子}——
都吸入\omterm{def:g7:what-is-respiration:respiration}{氧气}，释放\omterm{def:g7:what-is-respiration:respiration}{二氧化碳}。”

\textbf{9.} \omterm{def:g7:what-is-respiration:respiration}{氧气}下降，\omterm{prop:g7:what-is-respiration:shown}{澄清石灰水}变浑浊——
蘑菇在进行\omterm{def:g7:what-is-respiration:respiration}{呼吸作用}。它说明这一规律
同样适用于\omterm{def:g6:decomposers-and-soil:decomposers}{真菌}：\omterm{def:g6:decomposers-and-soil:decomposers}{真菌}既非\omterm{def:g1:plants-around-us:plant}{植物}，也非\omterm{def:g1:animals-around-us:animal}{动物}，没有例外。

\textbf{10.} 粮仓相当于仓库规模的 C 瓶：
潮湿保存的谷物开始\omterm{def:g2:from-seed-to-plant:germination}{萌发}，并剧烈地进行\omterm{def:g7:what-is-respiration:respiration}{呼吸作用}，
使储存的粮食升温、变质——管理员对干燥的
要求，正是对\omterm{def:g7:what-is-respiration:respiration}{呼吸作用}的控制。

\textbf{11.} 其中包括：池塘里的蜗牛和小\omterm{def:g1:animals-around-us:animal}{动物}、
淤泥中的\omterm{def:g6:decomposers-and-soil:decomposers}{分解者}、\omterm{def:g1:plants-around-us:plant}{植物}，以及绿色\omterm{def:g6:cell-unit-of-life:unicellular}{单细胞}\omterm{def:g6:exploring-our-environment:components}{生物}，
它们在黑暗时段同样需要\omterm{def:g7:what-is-respiration:respiration}{呼吸}。这一点很重要，因为池塘中的
所有\omterm{def:g6:exploring-our-environment:components}{生物}都依赖同一小份溶解氧储备——
这笔“\omterm{def:g7:what-is-respiration:respiration}{氧气}账”正是后面几章要展开的内容。

\textbf{12.} 只有 D——小鼠——显示出\omterm{def:g7:what-is-respiration:respiration}{呼吸}运动：
可见的胸廓起伏，服务于它的气体交换。每个
含有\omterm{def:g6:exploring-our-environment:components}{生物}的广口瓶都显示出\omterm{def:g7:what-is-respiration:respiration}{呼吸作用}：B 微弱，
C、D、F 明显，E 被白天的物质制造所掩盖，
但同样在进行\omterm{def:g7:what-is-respiration:respiration}{呼吸作用}。
\end{solution}
